\chapter{Les Faiseurs de Veuves}
\sessioninfo{Session 3}{Mardi 1\ier{} septembre 2026}

L'elfe les attendait devant une vaste demeure située à l'extérieur de Phandalin. Elle était là depuis cinq jours et, maintenant que les voyageurs avaient enfin atteint leur destination, elle les invita à entrer. Près de la porte, une étoile à huit branches semblable à un soleil attira l'attention de Morrÿs et d'Obélik. Tous deux reconnurent l'emblème des Faiseurs de Veuves, une organisation dont ils avaient déjà entendu le nom sans savoir précisément ce qui se cachait derrière.

La demeure était le quartier général de la guilde. Le groupe franchit l'entrée et gagna un petit patio baigné de lumière. L'elfe appela Yarno El Raïb, puis les nouveaux venus apprirent enfin le nom de celle qui les avait attendus si longtemps : Lyra.

Lorsqu'il les rejoignit, Yarno discutait avec un homme chauve nommé Destos. Il savait déjà qui étaient les nouveaux venus : Corrin Thinspike, l'étrange halfelin rencontré devant la prison de Waterdeep, l'avait averti de leur arrivée et les avait chaudement recommandés.

Kael, Morrÿs, Obélik et Alan lui racontèrent donc leur arrestation, leur évasion et les circonstances assez particulières qui les avaient conduits jusqu'à Phandalin. L'évocation de celle qu'ils appelaient désormais sans beaucoup d'affection « l'autre connasse » arracha à Yarno un petit rire amusé. Il voyait manifestement très bien de qui ils parlaient. D'après lui, elle arrêtait les gens à tour de bras, et les Faiseurs de Veuves cherchaient justement à rassembler suffisamment de preuves pour la faire tomber.

Cela ne réglait pourtant pas leur problème immédiat. Retourner à Waterdeep était impossible tant qu'ils y étaient recherchés. Rester à Phandalin ne l'était guère davantage : Yarno ne souhaitait pas attirer d'ennuis sur la guilde avant d'avoir réuni les preuves nécessaires.

Il allait donc falloir les envoyer ailleurs.

Yarno leur expliqua alors ce que faisaient réellement les Faiseurs de Veuves. Leur mission consistait notamment à retrouver les objets magiques susceptibles de devenir dangereux, à les mettre hors d'état de nuire et, lorsque cela s'avérait nécessaire, à les détruire. Ils n'étaient pas des collectionneurs et encore moins des pillards : leur rôle était de protéger ceux que de tels objets pouvaient menacer.

Une mission pourrait justement convenir au groupe. Elle les conduirait loin de Waterdeep et de ceux qui les recherchaient, sur les terres tropicales de Chult. L'endroit était sauvage, couvert d'une jungle dangereuse peuplée de créatures hostiles, venimeuses ou empoisonnées. Port Nyanzaru était la seule cité de l'île ; au-delà commençait un territoire qui promettait autant de périls que de découvertes.

L'un des sept princes marchands, un mage nommé Wakanga O'Tamu, avait fait appel aux Faiseurs de Veuves. Il recherchait un Gardien Bouclier ainsi que l'amulette permettant de le contrôler. Toute information permettant de retrouver l'un ou l'autre serait rémunérée ; rapporter le Gardien, l'amulette, ou les deux, vaudrait naturellement une prime autrement plus importante.

La guilde avait déjà deux vigies sur place : Iolka et Nera Vesper, deux sœurs que Yarno désignait simplement comme « la Paire ». Elles vivaient chez le commanditaire lorsqu'elles n'étaient pas en vadrouille. Il était donc possible que les aventuriers ne les trouvent pas dès leur arrivée.

Les Faiseurs de Veuves préféraient ne pas apparaître trop ouvertement dans cette affaire. Le groupe agirait en son nom propre et devrait se souvenir d'une chose essentielle : une fois à Chult, aucun d'entre eux ne ferait partie de la guilde. Ils seraient indépendants. Ils pourraient conserver les trésors abandonnés qu'ils découvriraient au cours de leurs recherches, mais il n'était pas question de voler ce qui appartenait à quelqu'un.

L'offre paraissait prometteuse. Yarno n'était cependant pas prêt à envoyer à l'autre bout de la Côte des Épées trois anciens prisonniers dont il ne savait presque rien sur la seule recommandation d'un halfelin capable de marcher au plafond.

Il avait prévu un test.

\scenebreak

Le groupe descendit au sous-sol, dans une salle d'armes aménagée pour l'entraînement. Chacun put s'équiper avant de rejoindre le terrain central. Lyra les y attendait. Yarno exposa une règle d'une simplicité rassurante : pour réussir l'épreuve, il leur suffisait de la toucher trois fois.

Puis il se tourna vers l'elfe et lui demanda d'y aller doucement.

Il ne fallait pas trop les abîmer.

Kael choisit de rester à distance et arma son arc. Morrÿs prépara son poignard, tandis que Destos, qui était clerc, se tenait prêt à lancer ses sortilèges et qu'Obélik s'avançait au contact. Lyra adopta une garde souple, les mains nues. L'objectif semblait modeste, mais les premiers assauts dissipèrent rapidement toute illusion : atteindre l'elfe serait déjà un exploit.

\illustration[0.8\textwidth]{essai-qg-fdv.png}
{\textit{Kael, Morrÿs, Obélik et Destos face à Lyra lors de l'épreuve.}}
{epreuve_lyra}

Lyra, elle, n'avait besoin ni d'arc ni d'arme. Lorsqu'elle parvenait au contact, un coup de poing ou la paume de sa main suffisait à faire disparaître brutalement toute l'énergie de sa cible. Un autre geste semblait parfois en rendre une partie, sans que ses adversaires comprennent vraiment ce qu'elle leur faisait.

Obélik fut le premier à subir ses coups, suivi de Destos, puis de Kael. Le goliath et le rôdeur finirent tous deux à terre. Aucun autre membre du groupe ne fut atteint. Malgré les coups reçus, Obélik réussit à toucher Lyra une première fois, puis à refermer ses bras autour d'elle et à l'immobiliser lors d'un nouvel assaut. Destos, resté à distance, parvint lui aussi à l'atteindre avec l'un de ses sortilèges.

Morrÿs tenta de participer à distance en lançant son poignard. Lyra dévia une première fois l'arme, qui alla se planter dans un mur. Le nain put la récupérer et tenta sa chance une seconde fois. Cette fois, l'elfe leva la jambe et détourna la lame d'un coup de pied. Le poignard fila au-dessus des combattants et se planta dans le plafond.

Bien trop haut pour Morrÿs.

Privé de son poignard, Morrÿs ne pouvait plus le récupérer pour tenter de le lancer une troisième fois. Yarno lui fit remarquer l'un des inconvénients des armes de jet : une fois lancées, encore fallait-il pouvoir les récupérer. La leçon aurait probablement gagné en dignité si le poignard n'était pas resté suspendu plusieurs mètres au-dessus de la tête du nain.

Au terme des échanges, le résultat était là. Entre le coup porté par Obélik, la prise avec laquelle il parvint à attraper Lyra et le sortilège de Destos, le groupe remplit l'objectif fixé par Yarno.

Yarno mit un terme à l'épreuve. Le résultat n'était peut-être pas élégant et plusieurs d'entre eux avaient passé une partie du combat au sol, mais il avait vu ce qu'il voulait voir. Ils avaient coopéré, couvert leurs faiblesses et continué jusqu'à remplir leur objectif. Cela suffisait pour leur confier la mission.

Il restait toutefois une dernière affaire à régler. Le poignard de Morrÿs était toujours planté dans le plafond.

Obélik sauta pour le décrocher et le rendit à son propriétaire.

\scenebreak

Yarno leur montra ensuite les chambres qui leur étaient réservées à l'étage. Elles donnaient sur une terrasse orientée plein sud, ouverte sur le jardin de la demeure. Phandalin n'était pas visible depuis là, mais, très loin, se dessinait la montagne dans laquelle ils avaient été emprisonnés.

Dans le jardin se dressaient quatre statues, une pour chacun des fondateurs des Faiseurs de Veuves. La guilde était née à Phandalin une trentaine d'années plus tôt et, depuis, ses membres portaient assistance aux populations jusque sur la Côte des Épées. Yarno leur confia également avoir été envoyé en prison, plusieurs années auparavant, par l'un de ces quatre fondateurs.

Il n'en dit pas davantage.

Les préparatifs du départ se précisèrent au fil de la soirée. Chacun recevrait cinquante pièces d'or, une potion de soin et une potion antipoison. Yarno leur remettrait également une carte de Chult avant leur départ, ainsi qu'une vieille photographie de leurs deux contacts.

Sur le cliché, Iolka et Nera Vesper posaient côte à côte sous un soleil éclatant. Les deux sœurs souriaient à l'objectif, suffisamment proches pour que le surnom choisi par Yarno paraisse parfaitement naturel. Si elles étaient présentes à Port Nyanzaru, le groupe saurait au moins qui chercher.

\illustration[0.8\textwidth]{La_Paire.png}
{\textit{Iolka et Nera Vesper, « la Paire ».}}
{la_paire}

Le voyage commencerait dès l'aube. Avec des vents favorables, il faudrait environ vingt et un jours de navigation pour atteindre Chult. Si les conditions se montraient mauvaises, la traversée pourrait en prendre près de quarante. Destos les accompagnerait : il avait participé à l'épreuve contre Lyra et ferait partie de l'expédition.

Yarno rappela une dernière fois que, là-bas, ils seraient seuls. Ils pourraient compter sur Iolka et Nera s'ils les trouvaient, mais ils ne porteraient ni le nom ni l'autorité des Faiseurs de Veuves. La guilde finançait la mission et leur fournissait un point de départ ; le reste leur appartiendrait.

\scenebreak

Une question demeurait : qu'allaient devenir Alan et Virgile ? Le tavernier ne pouvait guère retourner à Waterdeep avec les autres fugitifs. Une place lui fut proposée dans les cuisines de la demeure, proposition qui lui offrait à la fois un toit, du travail et un accès à des réserves qu'il ne tarda pas à examiner avec beaucoup d'intérêt.

Virgile prit une décision différente. Il ne voulait pas être considéré comme un déserteur et comptait retourner à Waterdeep. Son récit serait simple : il avait accompagné les prisonniers en raison des événements survenus à la prison, puis ceux-ci s'étaient joués de lui et avaient réussi à s'échapper.

Il préférait passer pour légèrement incompétent plutôt que pour un traître.

Avant son départ, son état nécessitait encore quelques soins. Lyra l'installa dans le patio et Destos lui prêta assistance grâce aux pouvoirs dont il disposait. Pendant ce temps, les autres explorèrent davantage la demeure et finirent par trouver les cuisines.

Alan y avait déjà pris ses marques.

Alan avait découvert une bonne bouteille dans les réserves et avait aussitôt décidé de l'utiliser pour préparer un chevreuil en sauce. Morrÿs s'intéressa au dessert et demanda s'il y aurait du chocolat. Alan ne connaissait pas cette denrée, mais la conversation donna une idée au groupe : s'ils trouvaient à Chult des épices intéressantes, ils lui en rapporteraient.

Le soir venu, tous se retrouvèrent dans le Salon Bleu pour partager le repas. Yarno profita du dîner pour leur parler davantage de leur destination. Port Nyanzaru était gouvernée par sept princes et princesses marchands. Leur titre ne désignait pas une lignée royale : ils étaient avant tout les personnes les plus riches et les plus influentes de la cité, et se partageaient ses responsabilités ainsi que différents domaines d'activité.

La chaleur serait un autre problème. Chult possédait un climat tropical auquel aucun d'entre eux n'était réellement préparé. Yarno observa notamment l'épaisse tenue de Destos, un homme venu du Nord habitué à se protéger du froid, et lui conseilla vivement d'abandonner quelques couches avant d'embarquer.

Le repas s'acheva, et chacun put enfin gagner sa chambre. Le lendemain, ils quitteraient Phandalin pour prendre la mer. Au bout du voyage les attendaient une cité inconnue, une jungle empoisonnée, deux vigies peut-être absentes et un Gardien Bouclier dont personne ne savait où il se trouvait.

Pour Alan, en revanche, l'avenir immédiat paraissait beaucoup plus simple.

Il avait trouvé une cuisine.
